% GENERATED FILE. Edit canonical lessons, then run npm run generate:books.
% canonical-source-hash: fnv1a64:00b528e3a8427ab7
% canonical-lessons: TE-C379-egabadu, TE-C379-padipovu, TE-C379-badhitudu, TE-C379-morapettuko, TE-C379-veduko

\chapter{Crowds and Cries}
\label{ch:te-crowds-and-cries}
\hlchaptermodality{379}

\begin{chapteropening}
\textbf{By the end of this chapter:} \emph{I can say to throng, to fall down, a victim, to petition and to beg.}

Complete the last lesson of chapter 379: I can say to throng, to fall down, a victim, to petition and to beg.
\end{chapteropening}

\section[\texorpdfstring{\te{ఎగబడు} (egabaḍu)}{ఎగబడు (egabaḍu)}]{\te{ఎగబడు} (egabaḍu) — to throng; to crowd in}

\label{lesson:TE-C379-egabadu}



\begin{quote}
Before the new one: say the Telugu for to grapple, then the Telugu for to gather round.
\end{quote}

\subsection*{You'll want to know: \te{ఎగబడు}}
\textbf{\te{ఎగబడు}} — \emph{egabaḍu} — \textquotedblleft{}to throng; to crowd in\textquotedblright{}.

\begin{quote}
\te{సినిమా} \te{టిక్కెట్ల} \te{కోసం} \te{జనం} \te{ఎగబడ్డారు}.
\end{quote}

\begin{quote}
\emph{sinimā ṭikkeṭla kōsaṁ janaṁ egabaḍḍāru.}
\end{quote}

\begin{quote}
People thronged for film tickets.
\end{quote}

Past (I …): \te{ఎగబడ్డాను}, \emph{egabaḍḍānu}.

It shows eagerness: everyone pushing to get there first.

\begin{cousinweb}
\textbf{\te{ఎగబడు}} (\emph{egabaḍu}) is a Telugu verb whose origin is not settled.
\end{cousinweb}

\begin{grammarlens}[title={What you've built}]
One more word for telling what happened.
\end{grammarlens}

\subsection*{Your turn}
\begin{itemize}
\raggedright
  \item \emph{Say it:} \emph{egabaḍu}
  \item \emph{Say it:} \emph{egabaḍu}, once more
  \item \emph{Say it:} \emph{sinimā ṭikkeṭla kōsaṁ janaṁ egabaḍḍāru.} — People thronged for film tickets.
  \item \emph{Recall:} say the Telugu for to grapple, then the Telugu for to gather round, then say \emph{egabaḍu} again
\end{itemize}

\begin{tcolorbox}[breakable,colback=teal!4,colframe=teal!35!black,title={Before you move on}]
What does \te{ఎగబడు} mean? (\textquotedblleft{}To throng; to crowd in\textquotedblright{}.) Say it once more.
\end{tcolorbox}
\section[\texorpdfstring{\te{పడిపోవు} (paḍipōvu)}{పడిపోవు (paḍipōvu)}]{\te{పడిపోవు} (paḍipōvu) — to fall down}

\label{lesson:TE-C379-padipovu}



\begin{quote}
Before the new one: say the Telugu for to gather round, then the Telugu for to throng.
\end{quote}

\subsection*{You'll want to know: \te{పడిపోవు}}
\textbf{\te{పడిపోవు}} — \emph{paḍipōvu} — \textquotedblleft{}to fall down\textquotedblright{}.

\begin{quote}
\te{తోపులాటలో} \te{ఒక} \te{ముసలాయన} \te{పడిపోయాడు}.
\end{quote}

\begin{quote}
\emph{tōpulāṭalō oka musalāyana paḍipōyāḍu.}
\end{quote}

\begin{quote}
An old man fell down in the crush.
\end{quote}

Past (I …): \te{పడిపోయాను}, \emph{paḍipōyānu}.

\te{పడు}, to fall, with \te{పోవు}: a fall that happens to you.

\begin{cousinweb}
\textbf{\te{పడిపోవు}} (\emph{paḍipōvu}) is a native compound: \emph{paḍu}, \textquotedblleft{}to fall\textquotedblright{} (DEDR 3852), and \emph{pōvu}, \textquotedblleft{}to go\textquotedblright{} (DEDR 4572).
\end{cousinweb}

\begin{grammarlens}[title={What you've built}]
One more word for telling what happened.
\end{grammarlens}

\subsection*{Your turn}
\begin{itemize}
\raggedright
  \item \emph{Say it:} \emph{paḍipōvu}
  \item \emph{Say it:} \emph{paḍipōvu}, once more
  \item \emph{Say it:} \emph{tōpulāṭalō oka musalāyana paḍipōyāḍu.} — An old man fell down in the crush.
  \item \emph{Recall:} say the Telugu for to gather round, then the Telugu for to throng, then say \emph{paḍipōvu} again
\end{itemize}

\begin{tcolorbox}[breakable,colback=teal!4,colframe=teal!35!black,title={Before you move on}]
What does \te{పడిపోవు} mean? (\textquotedblleft{}To fall down\textquotedblright{}.) Say it once more.
\end{tcolorbox}
\section[\texorpdfstring{\te{బాధితుడు} (bādhituḍu)}{బాధితుడు (bādhituḍu)}]{\te{బాధితుడు} (bādhituḍu) — a victim}

\label{lesson:TE-C379-badhitudu}



\begin{quote}
Before the new one: say the Telugu for to throng, then the Telugu for to fall down.
\end{quote}

\subsection*{You'll want to know: \te{బాధితుడు}}
\textbf{\te{బాధితుడు}} — \emph{bādhituḍu} — \textquotedblleft{}a victim\textquotedblright{}.

\begin{quote}
\te{బాధితులకు} \te{ప్రభుత్వం} \te{సహాయం} \te{చేసింది}.
\end{quote}

\begin{quote}
\emph{bādhitulaku prabhutvaṁ sahāyaṁ cēsindi.}
\end{quote}

\begin{quote}
The government helped the victims.
\end{quote}

Plural: \te{బాధితులు}, \emph{bādhitulu}.

You know \te{బాధ}, pain: \te{బాధితుడు}, one who is made to suffer.

\begin{cousinweb}
\textbf{\te{బాధితుడు}} (\emph{bādhituḍu}) is borrowed from Sanskrit \emph{bādhita}, \textquotedblleft{}hard-pressed, oppressed\textquotedblright{} (CDIAL 9205).
\end{cousinweb}

\begin{grammarlens}[title={What you've built}]
One more word for telling what happened.
\end{grammarlens}

\subsection*{Your turn}
\begin{itemize}
\raggedright
  \item \emph{Say it:} \emph{bādhituḍu}
  \item \emph{Say it:} \emph{bādhituḍu}, once more
  \item \emph{Say it:} \emph{bādhitulaku prabhutvaṁ sahāyaṁ cēsindi.} — The government helped the victims.
  \item \emph{Recall:} say the Telugu for to throng, then the Telugu for to fall down, then say \emph{bādhituḍu} again
\end{itemize}

\begin{tcolorbox}[breakable,colback=teal!4,colframe=teal!35!black,title={Before you move on}]
What does \te{బాధితుడు} mean? (\textquotedblleft{}A victim\textquotedblright{}.) Say it once more.
\end{tcolorbox}
\section[\texorpdfstring{\te{మొరపెట్టుకో} (morapeṭṭukō)}{మొరపెట్టుకో (morapeṭṭukō)}]{\te{మొరపెట్టుకో} (morapeṭṭukō) — to petition; to plead}

\label{lesson:TE-C379-morapettuko}



\begin{quote}
Before the new one: say the Telugu for to fall down, then the Telugu for a victim.
\end{quote}

\subsection*{You'll want to know: \te{మొరపెట్టుకో}}
\textbf{\te{మొరపెట్టుకో}} — \emph{morapeṭṭukō} — \textquotedblleft{}to petition; to plead\textquotedblright{}.

\begin{quote}
\te{బాధితులు} \te{అధికారులకు} \te{మొరపెట్టుకున్నారు}.
\end{quote}

\begin{quote}
\emph{bādhitulu adhikārulaku morapeṭṭukunnāru.}
\end{quote}

\begin{quote}
The victims pleaded with the officials.
\end{quote}

Past (I …): \te{మొరపెట్టుకున్నాను}, \emph{morapeṭṭukunnānu}.

To whom takes -\te{కు}: \te{అధికారులకు} \te{మొరపెట్టుకో}.

\begin{cousinweb}
\textbf{\te{మొరపెట్టుకో}} (\emph{morapeṭṭukō}) is a native compound: \emph{moṟa}, \textquotedblleft{}a cry, a scream\textquotedblright{} (DEDR 5013), and \emph{peṭṭukō}, \textquotedblleft{}to put for oneself\textquotedblright{} (DEDR 4390).
\end{cousinweb}

\begin{grammarlens}[title={What you've built}]
One more word for telling what happened.
\end{grammarlens}

\subsection*{Your turn}
\begin{itemize}
\raggedright
  \item \emph{Say it:} \emph{morapeṭṭukō}
  \item \emph{Say it:} \emph{morapeṭṭukō}, once more
  \item \emph{Say it:} \emph{bādhitulu adhikārulaku morapeṭṭukunnāru.} — The victims pleaded with the officials.
  \item \emph{Recall:} say the Telugu for to fall down, then the Telugu for a victim, then say \emph{morapeṭṭukō} again
\end{itemize}

\begin{tcolorbox}[breakable,colback=teal!4,colframe=teal!35!black,title={Before you move on}]
What does \te{మొరపెట్టుకో} mean? (\textquotedblleft{}To petition; to plead\textquotedblright{}.) Say it once more.
\end{tcolorbox}
\section[\texorpdfstring{\te{వేడుకో} (vēḍukō)}{వేడుకో (vēḍukō)}]{\te{వేడుకో} (vēḍukō) — to beg; to entreat}

\label{lesson:TE-C379-veduko}



\begin{quote}
Before the new one: say the Telugu for a victim, then the Telugu for to petition.
\end{quote}

\subsection*{You'll want to know: \te{వేడుకో}}
\textbf{\te{వేడుకో}} — \emph{vēḍukō} — \textquotedblleft{}to beg; to entreat\textquotedblright{}.

\begin{quote}
\te{ఆమె} \te{దొంగను} \te{తన} \te{సంచి} \te{ఇవ్వమని} \te{వేడుకుంది}.
\end{quote}

\begin{quote}
\emph{āme doṅganu tana sañci ivvamani vēḍukundi.}
\end{quote}

\begin{quote}
She begged the thief to give her bag back.
\end{quote}

Past (I …): \te{వేడుకున్నాను}, \emph{vēḍukunnānu}.

Gentler than \te{బతిమాలు}; in the old stories, the weak beg the strong.

\begin{cousinweb}
\textbf{\te{వేడుకో}} (\emph{vēḍukō}) is native Dravidian: DEDR 5528 lists Telugu \emph{vēḍu}, \textquotedblleft{}to pray, beg, entreat\textquotedblright{}, here with the self-ending \emph{-kō}.
\end{cousinweb}

\begin{grammarlens}[title={What you've built}]
One more word for telling what happened.
\end{grammarlens}

\subsection*{Your turn}
\begin{itemize}
\raggedright
  \item \emph{Say it:} \emph{vēḍukō}
  \item \emph{Say it:} \emph{vēḍukō}, once more
  \item \emph{Say it:} \emph{āme doṅganu tana sañci ivvamani vēḍukundi.} — She begged the thief to give her bag back.
  \item \emph{Recall:} say the Telugu for a victim, then the Telugu for to petition, then say \emph{vēḍukō} again
\end{itemize}

\begin{tcolorbox}[breakable,colback=teal!4,colframe=teal!35!black,title={Before you move on}]
What does \te{వేడుకో} mean? (\textquotedblleft{}To beg; to entreat\textquotedblright{}.) Say it once more.
\end{tcolorbox}
